\documentclass[11pt]{amsart}

\usepackage[T1]{fontenc}
\usepackage{lmodern}
\usepackage{microtype}
\usepackage{amsmath,amssymb,amsthm}
\usepackage[hidelinks]{hyperref}

\hypersetup{
  pdftitle={Kite families with Theta(n^3) failure of the Zhou-Wang-Chai distance-Laplacian bound},
  pdfkeywords={distance Laplacian matrix, Wiener index, eigenvalue sum, counterexample, unicyclic graph}
}

\newtheorem{theorem}{Theorem}
\newtheorem{corollary}[theorem]{Corollary}
\theoremstyle{remark}
\newtheorem{remark}[theorem]{Remark}

\newcommand{\DL}{D^{L}}
\newcommand{\one}{\mathbf 1}
\newcommand{\gap}{\Delta}

\title[Cubic failure of the Zhou--Wang--Chai bound]{Kite Families with $\Theta(n^3)$ Failure of the Zhou--Wang--Chai Distance-Laplacian Bound}
\date{}

\subjclass[2020]{05C50, 05C12, 15A18}
\keywords{distance Laplacian matrix, Wiener index, eigenvalue sum, counterexample, unicyclic graph}

\begin{document}

\begin{abstract}
Let $\partial_1^L(G)\ge\cdots\ge\partial_n^L(G)=0$ be the distance Laplacian eigenvalues of a connected graph $G$, and write
$U_r(G)=\sum_{i=1}^r\partial_i^L(G)$.
Zhou, Wang, and Chai proposed the bound
\[
U_r(G)\le W(G)+\binom{r+2}{3}.
\]
Huang \cite{Huang} has already shown that this bound fails for $K_{1,3}$ at $r=2$ and for $K_{1,4}$ at $r=3$, so the unrestricted conjecture is false; our point is that the failure is neither sporadic nor small. We show that failure persists throughout an infinite unicyclic family. If $H_n$ is the kite graph obtained from the path $v_1v_2\cdots v_n$ by adding $v_1v_3$, then every $H_n$, $n\ge5$, violates the bound at $r=n-2$, with gap at least $(3n-14)/2$. Along $n=3m$, a second choice $r=2m$ gives a violation of order $n^3$. The graph $H_5$ has minimum possible order among unicyclic counterexamples.
\end{abstract}

\maketitle

\section{The counterexamples}

For a connected graph $G$, let $D(G)=(d_G(u,v))$ be its distance matrix, let
\[
\tau_G(v)=\sum_{u\in V(G)}d_G(u,v),
\qquad
W(G)=\frac12\sum_{v\in V(G)}\tau_G(v),
\]
and let
\[
\DL(G)=\operatorname{diag}(\tau_G(v):v\in V(G))-D(G)
\]
be its distance Laplacian matrix \cite{AouchicheHansen}.
Write
\[
\partial_1^L(G)\ge\cdots\ge\partial_n^L(G)=0,
\qquad
U_r(G)=\sum_{i=1}^r\partial_i^L(G),
\]
and define the excess
\[
\gap_r(G)=U_r(G)-W(G)-\binom{r+2}{3}.
\]
Zhou, Wang, and Chai \cite{ZhouWangChai} conjectured that $\gap_r(G)\le0$ for every connected graph and every $2\le r\le n$. A recent preprint of Huang \cite{Huang} classifies the diameter-two case and finds precisely the exceptional pairs $(K_{1,3},2)$ and $(K_{1,4},3)$. Our contribution is an infinite family inside the more restrictive class of unicyclic graphs.

For $n\ge4$, let $H_n$ have vertex set $\{v_1,\ldots,v_n\}$ and edge set
\[
E(H_n)=\{v_iv_{i+1}:1\le i<n\}\cup\{v_1v_3\}.
\]
Thus $H_n$ is the standard kite graph $\operatorname{Ki}_{n,3}$: a triangle with a pendant path \cite{LinDuZhou}.

\begin{theorem}\label{thm:main}
For every $n\ge5$,
\[
\gap_{n-2}(H_n)\ge\frac{3n-14}{2}>0.
\]
Moreover, if $n=3m$ with $m\ge4$, then
\[
\gap_{2m}(H_{3m})
\ge
\frac{5m^3-12m^2-17m+12}{6}>0.
\]
In particular, the violation is $\Theta(n^3)$ along the subfamily $H_{3m}$.
\end{theorem}

\begin{proof}
For every $x\in\mathbb R^{V(G)}$,
\begin{equation}\label{eq:quadratic}
x^{\mathsf T}\DL(G)x
=\sum_{\{u,v\}\subseteq V(G)}d_G(u,v)(x_u-x_v)^2.
\end{equation}
Hence $\DL(G)$ is positive semidefinite and, since $G$ is connected,
$\ker\DL(G)=\operatorname{span}\{\one\}$.
Therefore
\begin{equation}\label{eq:minmax}
\partial_{n-1}^L(G)
=
\min_{\substack{x\ne0\\x\perp\one}}
\frac{x^{\mathsf T}\DL(G)x}{x^{\mathsf T}x}.
\end{equation}

Only the distances $d(v_1,v_j)$, $3\le j\le n$, decrease when $v_1v_3$ is added to $P_n$, and each decreases by one. Thus
\begin{equation}\label{eq:wiener}
W(H_n)=W(P_n)-(n-2)=\binom{n+1}{3}-(n-2).
\end{equation}
Take $x=e_3-e_4$. The relevant transmissions are
\[
\tau_{H_n}(v_3)=2+\frac{(n-3)(n-2)}2
=\frac{n^2-5n+10}{2}
\]
and
\[
\tau_{H_n}(v_4)=5+\frac{(n-4)(n-3)}2
=\frac{n^2-7n+22}{2}.
\]
Since $x\perp\one$ and $d(v_3,v_4)=1$, equations \eqref{eq:quadratic}--\eqref{eq:minmax} give
\[
\partial_{n-1}^L(H_n)
\le
\frac{\tau_{H_n}(v_3)+\tau_{H_n}(v_4)+2}{2}
=
\frac{n^2-6n+18}{2}.
\]
Also $\operatorname{tr}\DL(H_n)=2W(H_n)$, so
\[
U_{n-2}(H_n)=2W(H_n)-\partial_{n-1}^L(H_n).
\]
Using \eqref{eq:wiener},
\begin{align*}
\gap_{n-2}(H_n)
&\ge
W(H_n)-\binom n3-\frac{n^2-6n+18}{2}\\
&=
\frac{3n-14}{2}.
\end{align*}
This proves the first assertion.

Now let $n=3m$ and
\[
S=\{v_1,\ldots,v_m\}\cup\{v_{2m+1},\ldots,v_{3m}\}.
\]
Ky Fan's variational principle \cite{Fan} gives
\begin{equation}\label{eq:kyfan}
U_{2m}(H_{3m})
=
\max_{Q^{\mathsf T}Q=I_{2m}}
\operatorname{tr}\!\left(Q^{\mathsf T}\DL(H_{3m})Q\right)
\ge
\sum_{v\in S}\tau_{H_{3m}}(v),
\end{equation}
where $Q$ ranges over real $3m\times2m$ matrices. For the path $P_{3m}$,
\begin{align*}
\sum_{v\in S}\tau_{P_{3m}}(v)
&=
2\sum_{i=1}^{m}
\left[\binom{i}{2}+\binom{3m-i+1}{2}\right]\\
&=
\frac{20m^3-2m}{3}.
\end{align*}
Passing from $P_{3m}$ to $H_{3m}$ shortens the $3m-2$ pairs
$\{v_1,v_j\}$, $3\le j\le3m$. Every such pair contributes once to the selected transmission sum through $v_1\in S$, and $2m-2$ of them contribute a second time because $v_j\in S$. Hence
\begin{equation}\label{eq:selected}
\sum_{v\in S}\tau_{H_{3m}}(v)
=
\frac{20m^3-2m}{3}-(5m-4)
=
\frac{20m^3-17m+12}{3}.
\end{equation}
Combining \eqref{eq:wiener}, \eqref{eq:kyfan}, and \eqref{eq:selected} yields
\[
\gap_{2m}(H_{3m})
\ge
\frac{5m^3-12m^2-17m+12}{6}.
\]
For $p(m)=5m^3-12m^2-17m+12$, one has $p(4)=72$ and
\[
p(m+1)-p(m)=3(m+1)(5m-8)>0
\qquad(m\ge4),
\]
so the lower bound is positive. It is $\frac{5}{162}n^3+O(n^2)$, whereas
$\gap_{2m}(H_{3m})\le U_{2m}(H_{3m})\le2W(H_{3m})=O(n^3)$.
Thus the violation is $\Theta(n^3)$.
\end{proof}

\begin{remark}
The unicyclic restriction is what separates Theorem~\ref{thm:main} from the unrestricted statement; it is not what makes the failure infinite, nor what makes it large. Under the unrestricted formulation, the same argument with $x=e_2-e_3$ gives
\[
\gap_{n-2}(P_n)\ge\frac{3n-10}{2}>0
\qquad(n\ge4).
\]
Moreover, the sum $\sum_{v\in S}\tau_{P_{3m}}(v)$ evaluated in the proof of Theorem~\ref{thm:main} already exceeds $W(P_{3m})+\binom{2m+2}{3}$ for every $m\ge3$, so by Ky Fan's principle \cite{Fan} the paths $P_{3m}$ by themselves violate the bound by $\Theta(n^3)$. Thus paths already form an infinite family with cubic excess; Theorem~\ref{thm:main} shows that failure of both orders persists after imposing unicyclicity.
\end{remark}

\begin{corollary}\label{cor:five}
The graph $H_5$ has minimum possible order among unicyclic counterexamples. Its distance Laplacian matrix is
\[
\DL(H_5)=
\begin{pmatrix}
 7&-1&-1&-2&-3\\
-1& 7&-1&-2&-3\\
-1&-1& 5&-1&-2\\
-2&-2&-1& 6&-1\\
-3&-3&-2&-1& 9
\end{pmatrix},
\]
and
\[
\det(xI-\DL(H_5))
=x(x-6)(x-8)(x^2-20x+95).
\]
Consequently,
\[
\operatorname{Spec}(\DL(H_5))
=\{10+\sqrt5,\,8,\,10-\sqrt5,\,6,\,0\},
\]
so $W(H_5)=17$ and
\[
U_3(H_5)=28>27=17+\binom53.
\]
\end{corollary}

\begin{proof}
The displayed matrix and characteristic polynomial follow directly from the distances in $H_5$. The only connected unicyclic graphs on fewer than five vertices are $C_3$, $C_4$, and $H_4$, whose distance Laplacian spectra are, respectively,
\[
\{3,3,0\},
\qquad
\{6,6,4,0\},
\qquad
\{7,5,4,0\}.
\]
These spectra verify the conjectured inequality for every relevant $r$, while the calculation above shows that $H_5$ violates it at $r=3$.
\end{proof}

\begin{thebibliography}{9}

\bibitem{AouchicheHansen}
M.~Aouchiche and P.~Hansen,
\emph{Two Laplacians for the distance matrix of a graph},
Linear Algebra Appl. \textbf{439} (2013), no.~1, 21--33.
\href{https://doi.org/10.1016/j.laa.2013.02.030}{doi:10.1016/j.laa.2013.02.030}.

\bibitem{Fan}
K.~Fan,
\emph{On a theorem of Weyl concerning eigenvalues of linear transformations I},
Proc. Natl. Acad. Sci. USA \textbf{35} (1949), no.~11, 652--655.
\href{https://doi.org/10.1073/pnas.35.11.652}{doi:10.1073/pnas.35.11.652}.

\bibitem{Huang}
S.~Huang,
\emph{On a distance Laplacian analog of Brouwer's conjecture for several classes of graphs},
arXiv:2606.06945 [math.CO], 2026.

\bibitem{LinDuZhou}
H.~Lin, Z.~Du, and B.~Zhou,
\emph{On the first two largest distance Laplacian eigenvalues of unicyclic graphs},
Linear Algebra Appl. \textbf{546} (2018), 289--307.
\href{https://doi.org/10.1016/j.laa.2018.02.013}{doi:10.1016/j.laa.2018.02.013}.

\bibitem{ZhouWangChai}
Y.~Zhou, L.~Wang, and Y.~Chai,
\emph{Brouwer type conjecture for the eigenvalues of distance Laplacian matrix of a graph},
Comput. Appl. Math. \textbf{44} (2025), Paper No.~138.
\href{https://doi.org/10.1007/s40314-025-03095-0}{doi:10.1007/s40314-025-03095-0}.

\end{thebibliography}

\end{document}